\documentclass[10pt,a4paper]{amsart}
\usepackage{amsmath,amssymb,amsthm,graphicx}
% Retained preamble for the Stacks Project chapter morphisms.tex (Morphisms of Schemes).
% The chapter file pinned-original-source.tex carries only \input{preamble} on its first line, so no
% preamble text lives inside the chapter itself.  The Stacks Project uses one single shared preamble for
% every chapter; the snapshot stages it verbatim in sources/stacks-chow/preamble.tex.  The block below is
% copied mechanically, line for line, from that file's lines 14-18 (the Xy-pic setup) and lines 154-217
% (the native theorem declarations together with the mathematical macros).  Nothing was added, removed or
% redefined.  The wrapper class and the amsmath/amssymb/amsthm/graphicx base packages are supplied by the
% assembler, so the class/\externaldocument part of the shared preamble is not repeated here.
% For commutative diagrams we use Xy-pic
\usepackage[all]{xy}

% We use 2cell for 2-commutative diagrams.
\xyoption{2cell}
% Theorem environments.
%
\theoremstyle{plain}
\newtheorem{theorem}[subsection]{Theorem}
\newtheorem{proposition}[subsection]{Proposition}
\newtheorem{lemma}[subsection]{Lemma}

\theoremstyle{definition}
\newtheorem{definition}[subsection]{Definition}
\newtheorem{example}[subsection]{Example}
\newtheorem{exercise}[subsection]{Exercise}
\newtheorem{situation}[subsection]{Situation}

\theoremstyle{remark}
\newtheorem{remark}[subsection]{Remark}
\newtheorem{remarks}[subsection]{Remarks}

\numberwithin{equation}{subsection}

% Macros
%
\def\lim{\mathop{\mathrm{lim}}\nolimits}
\def\colim{\mathop{\mathrm{colim}}\nolimits}
\def\Spec{\mathop{\mathrm{Spec}}}
\def\Hom{\mathop{\mathrm{Hom}}\nolimits}
\def\Ext{\mathop{\mathrm{Ext}}\nolimits}
\def\SheafHom{\mathop{\mathcal{H}\!\mathit{om}}\nolimits}
\def\SheafExt{\mathop{\mathcal{E}\!\mathit{xt}}\nolimits}
\def\Sch{\mathit{Sch}}
\def\Mor{\mathop{\mathrm{Mor}}\nolimits}
\def\Ob{\mathop{\mathrm{Ob}}\nolimits}
\def\Sh{\mathop{\mathit{Sh}}\nolimits}
\def\NL{\mathop{N\!L}\nolimits}
\def\CH{\mathop{\mathrm{CH}}\nolimits}
\def\proetale{{pro\text{-}\acute{e}tale}}
\def\etale{{\acute{e}tale}}
\def\QCoh{\mathit{QCoh}}
\def\Ker{\mathop{\mathrm{Ker}}}
\def\Im{\mathop{\mathrm{Im}}}
\def\Coker{\mathop{\mathrm{Coker}}}
\def\Coim{\mathop{\mathrm{Coim}}}

% Boxtimes
%
\DeclareMathSymbol{\boxtimes}{\mathbin}{AMSa}{"02}

%
% Macros for moduli stacks/spaces
%
\def\QCohstack{\mathcal{QC}\!\mathit{oh}}
\def\Cohstack{\mathcal{C}\!\mathit{oh}}
\def\Spacesstack{\mathcal{S}\!\mathit{paces}}
\def\Quotfunctor{\mathrm{Quot}}
\def\Hilbfunctor{\mathrm{Hilb}}
\def\Curvesstack{\mathcal{C}\!\mathit{urves}}
\def\Polarizedstack{\mathcal{P}\!\mathit{olarized}}
\def\Complexesstack{\mathcal{C}\!\mathit{omplexes}}
% \Pic is the operator that assigns to X its picard group, usage \Pic(X)
% \Picardstack_{X/B} denotes the Picard stack of X over B
% \Picardfunctor_{X/B} denotes the Picard functor of X over B
\def\Pic{\mathop{\mathrm{Pic}}\nolimits}
\def\Picardstack{\mathcal{P}\!\mathit{ic}}
\def\Picardfunctor{\mathrm{Pic}}
\def\Deformationcategory{\mathcal{D}\!\mathit{ef}}

% --- appended 2026-10-12 (Chief): the project own definitions of the reference/slogan/history
% display environments, copied mechanically from the shared Stacks preamble lines 8-12
% (\usepackage{verbatim} plus the three \newenvironment lines, which the chapter preamble keeps
% verbatim from the project preamble). Nothing added, removed or redefined: these environments
% carry no mathematical content and the project itself typesets them as comments. Without these
% three lines a counted statement whose own text uses them fails with Environment slogan undefined.
\usepackage{verbatim}
\newenvironment{reference}{\comment}{\endcomment}
\newenvironment{slogan}{\comment}{\endcomment}
\newenvironment{history}{\comment}{\endcomment}
\usepackage{cleveref}
\crefname{originallemma}{original lemma}{original lemmas}
\Crefname{originallemma}{Original lemma}{Original lemmas}
\makeatletter
\newlabel{lemma-finite-type-ample-very-ample}{{01VU}{source}{}{sourcepointer-1}{}}
\newlabel{lemma-finite-type-ample-very-ample@cref}{{[originallemma][0][]01VU}{[1][0][]source}{}{}{}}
\makeatother
\crefname{originallemma}{original lemma}{original lemmas}
\Crefname{originallemma}{Original lemma}{Original lemmas}
\makeatletter
\newlabel{lemma-immersion-permanence}{{07RK}{source}{}{sourcepointer-2}{}}
\newlabel{lemma-immersion-permanence@cref}{{[originallemma][0][]07RK}{[1][0][]source}{}{}{}}
\makeatother
\crefname{originallemma}{original lemma}{original lemmas}
\Crefname{originallemma}{Original lemma}{Original lemmas}
\makeatletter
\newlabel{lemma-quasi-compact-immersion}{{01RG}{source}{}{sourcepointer-3}{}}
\newlabel{lemma-quasi-compact-immersion@cref}{{[originallemma][0][]01RG}{[1][0][]source}{}{}{}}
\makeatother
\crefname{originallemma}{original lemma}{original lemmas}
\Crefname{originallemma}{Original lemma}{Original lemmas}
\makeatletter
\newlabel{lemma-quasi-projective-properties}{{01VX}{source}{}{sourcepointer-4}{}}
\newlabel{lemma-quasi-projective-properties@cref}{{[originallemma][0][]01VX}{[1][0][]source}{}{}{}}
\makeatother
\crefname{originallemma}{original lemma}{original lemmas}
\Crefname{originallemma}{Original lemma}{Original lemmas}
\makeatletter
\newlabel{lemma-very-ample-finite-type-over-affine}{{02NP}{source}{}{sourcepointer-5}{}}
\newlabel{lemma-very-ample-finite-type-over-affine@cref}{{[originallemma][0][]02NP}{[1][0][]source}{}{}{}}
\makeatother
\crefname{originallemma}{original lemma}{original lemmas}
\Crefname{originallemma}{Original lemma}{Original lemmas}
\makeatletter
\newlabel{properties-lemma-directed-colimit-finite-presentation}{{01PK}{source}{}{sourcepointer-6}{}}
\newlabel{properties-lemma-directed-colimit-finite-presentation@cref}{{[originallemma][0][]01PK}{[1][0][]source}{}{}{}}
\makeatother
\crefname{originallemma}{original lemma}{original lemmas}
\Crefname{originallemma}{Original lemma}{Original lemmas}
\makeatletter
\newlabel{schemes-lemma-push-forward-quasi-coherent}{{01LC}{source}{}{sourcepointer-7}{}}
\newlabel{schemes-lemma-push-forward-quasi-coherent@cref}{{[originallemma][0][]01LC}{[1][0][]source}{}{}{}}
\makeatother
\crefname{originallemma}{original lemma}{original lemmas}
\Crefname{originallemma}{Original lemma}{Original lemmas}
\makeatletter
\newlabel{schemes-lemma-quasi-compact-permanence}{{03GI}{source}{}{sourcepointer-8}{}}
\newlabel{schemes-lemma-quasi-compact-permanence@cref}{{[originallemma][0][]03GI}{[1][0][]source}{}{}{}}
\makeatother
\title{A quasi-projective morphism over a quasi-compact quasi-separated base is an open immersion into a projective morphism}
\author{The Stacks Project Authors}
\date{stacks-project:morphisms.tex:local-snapshot-20261012; GNU FDL 1.2 or later}
\begin{document}
\maketitle
\noindent\textit{Editorial note.} All delivered mathematical text is the author's own, copied from the pinned Stacks Project chapter file pinned-original-source.tex; no statement or proof step was rewritten, reordered or removed.\par\medskip
\noindent\textit{Editorial note.} Following the multi-claim convention for this source, source\_url is the source-level URL https://stacks.math.columbia.edu/ and the source label lemma-quasi-projective-open-projective is carried in claim\_key and here rather than in a per-claim tag URL.\par\medskip
\noindent\textit{Editorial note.} The preamble is the chapter's shared Stacks Project preamble, retained as sources/stacks-morphisms/retained-preamble.tex: a line-for-line excerpt of the shared Stacks preamble staged with this snapshot. No macro was invented or redefined, and this card adds no preamble text of its own.\par\medskip
\noindent\textit{Editorial note.} Statement span 10717-10722 and proof span 10724-10774 come from the halfA pool card operations/pool-cards/stacks-morphisms-halfA.json and were re-read in the pinned chapter file before this card was written.\par\medskip
\noindent\textit{Editorial note.} Re-checked by hand before carding: this candidate was caught by the lane's thin-proof prefilter because its last sentence contains the phrase 'by definition'. Reading the original shows the proof is a full construction (very ample twist, colimit approximation of the pushforward, lifting of finitely many sections, immersion check, and the quasi-compact immersion step), so the prefilter hit was incidental; the lemma is delivered as the author's own complete argument, and the prefilter reason is recorded here per the fleet rule of 2026-10-14 that a prefilter hit on a proof with at least 800 bytes or 100 lines must be re-read by hand.\par\medskip
\noindent\textit{Editorial note.} 8 cross-references occur inside the counted spans; they are listed in external\_references as navigation pointers only.\par\medskip
\noindent\textit{Allowed standard prerequisites (editorial).}\begin{itemize}
\item For a finite type morphism over a quasi-compact base a sufficiently high power of an ample sheaf is very ample, cited from the same chapter. Cited at proof lines 10724-10727. Lets the proof replace the given sheaf by a power that is very ample.
\item The pushforward of the structure-like sheaf along a quasi-projective morphism is quasi-coherent, and quasi-projective morphisms are quasi-compact and separated, cited from the schemes chapter and the same chapter. Cited at proof lines 10730-10734. Makes the pushforward of the sheaf a quasi-coherent module and gives the quasi-compactness used later.
\item A quasi-coherent module is the filtered colimit of its finite type submodules, cited from the properties chapter. Cited at proof lines 10735-10738. Produces the finite type modules over which the construction of the morphism is made.
\item The description of very ample sheaves on an affine chart by finitely many generating sections and the resulting immersion into a projective space, and permanence of immersions. Cited at proof lines 10744-10749 and 10766-10767. Give the local immersions and show that the globally constructed morphism is an immersion.
\item The statement that an immersion which is quasi-compact factors through an open immersion into a closed subscheme of its target, cited from the same chapter. Cited at proof lines 10771-10773. Produces the intermediate scheme X' inside the projective bundle.
\end{itemize}\medskip
\ifdefined\hypertarget\hypertarget{sourcepointer-1}{}\fi
\noindent\textit{Source reference.} Original lemma 01VU: focused-seeds/AGP/sources/stacks-snapshot-20261011/tags.txt, source lines 2459--2459. Same-chapter reference used in the retained author text at source line(s) 9056, 10727. It is supplied only as a navigation pointer so that the retained text compiles; the counted item delivers the author's own statement and the author's complete proof, and no part of the referenced item is copied in or substituted for the counted argument.\par\medskip
\ifdefined\hypertarget\hypertarget{sourcepointer-2}{}\fi
\noindent\textit{Source reference.} Original lemma 07RK: focused-seeds/AGP/sources/stacks-snapshot-20261011/tags.txt, source lines 9870--9870. Same-chapter reference used in the retained author text at source line(s) 361, 9510, 9691, 10764. It is supplied only as a navigation pointer so that the retained text compiles; the counted item delivers the author's own statement and the author's complete proof, and no part of the referenced item is copied in or substituted for the counted argument.\par\medskip
\ifdefined\hypertarget\hypertarget{sourcepointer-3}{}\fi
\noindent\textit{Source reference.} Original lemma 01RG: focused-seeds/AGP/sources/stacks-snapshot-20261011/tags.txt, source lines 2301--2301. Same-chapter reference used in the retained author text at source line(s) 1185, 1204, 2351, 10712, 10770. It is supplied only as a navigation pointer so that the retained text compiles; the counted item delivers the author's own statement and the author's complete proof, and no part of the referenced item is copied in or substituted for the counted argument.\par\medskip
\ifdefined\hypertarget\hypertarget{sourcepointer-4}{}\fi
\noindent\textit{Source reference.} Original lemma 01VX: focused-seeds/AGP/sources/stacks-snapshot-20261011/tags.txt, source lines 2462--2462. Same-chapter reference used in the retained author text at source line(s) 10733. It is supplied only as a navigation pointer so that the retained text compiles; the counted item delivers the author's own statement and the author's complete proof, and no part of the referenced item is copied in or substituted for the counted argument.\par\medskip
\ifdefined\hypertarget\hypertarget{sourcepointer-5}{}\fi
\noindent\textit{Source reference.} Original lemma 02NP: focused-seeds/AGP/sources/stacks-snapshot-20261011/tags.txt, source lines 3488--3488. Same-chapter reference used in the retained author text at source line(s) 9604, 10750. It is supplied only as a navigation pointer so that the retained text compiles; the counted item delivers the author's own statement and the author's complete proof, and no part of the referenced item is copied in or substituted for the counted argument.\par\medskip
\ifdefined\hypertarget\hypertarget{sourcepointer-6}{}\fi
\noindent\textit{Source reference.} Original lemma 01PK: focused-seeds/AGP/sources/stacks-snapshot-20261011/tags.txt, source lines 2233--2233. Same-chapter reference used in the retained author text at source line(s) 10734. It is supplied only as a navigation pointer so that the retained text compiles; the counted item delivers the author's own statement and the author's complete proof, and no part of the referenced item is copied in or substituted for the counted argument.\par\medskip
\ifdefined\hypertarget\hypertarget{sourcepointer-7}{}\fi
\noindent\textit{Source reference.} Original lemma 01LC: focused-seeds/AGP/sources/stacks-snapshot-20261011/tags.txt, source lines 2081--2081. Same-chapter reference used in the retained author text at source line(s) 262, 404, 882, 1765, 1924, 2256, 2296, 8726, 9175, 10731, 13903. It is supplied only as a navigation pointer so that the retained text compiles; the counted item delivers the author's own statement and the author's complete proof, and no part of the referenced item is copied in or substituted for the counted argument.\par\medskip
\ifdefined\hypertarget\hypertarget{sourcepointer-8}{}\fi
\noindent\textit{Source reference.} Original lemma 03GI: focused-seeds/AGP/sources/stacks-snapshot-20261011/tags.txt, source lines 4529--4529. Same-chapter reference used in the retained author text at source line(s) 1003, 2371, 4298, 8937, 8952, 9829, 10100, 10769, 13984, 13991, 14143. It is supplied only as a navigation pointer so that the retained text compiles; the counted item delivers the author's own statement and the author's complete proof, and no part of the referenced item is copied in or substituted for the counted argument.\par\medskip
\begin{lemma}
\label{lemma-quasi-projective-open-projective}
Let $f : X \to S$ be a quasi-projective morphism with $S$ quasi-compact
and quasi-separated. Then $f$ factors as $X \to X' \to S$ where $X \to X'$
is an open immersion and $X' \to S$ is projective.
\end{lemma}

\begin{proof}
Let $\mathcal{L}$ be $f$-ample. Since $f$ is of finite type and $S$ is
quasi-compact $\mathcal{L}^{\otimes n}$ is $f$-very ample for some $n > 0$, see
Lemma \ref{lemma-finite-type-ample-very-ample}.
Replace $\mathcal{L}$ by $\mathcal{L}^{\otimes n}$.
Write $\mathcal{F} = f_*\mathcal{L}$. This is a quasi-coherent
$\mathcal{O}_S$-module by 
Schemes, Lemma \ref{schemes-lemma-push-forward-quasi-coherent}
(quasi-projective morphisms are quasi-compact and
separated, see Lemma \ref{lemma-quasi-projective-properties}). By
Properties, Lemma \ref{properties-lemma-directed-colimit-finite-presentation}
we can find a directed set $I$ and a system of
finite type quasi-coherent $\mathcal{O}_S$-modules $\mathcal{E}_i$
over $I$ such that $\mathcal{F} = \colim \mathcal{E}_i$.
Consider the compositions
$\psi_i : f^*\mathcal{E}_i \to f^*\mathcal{F} \to \mathcal{L}$.
Choose a finite affine open covering $S = \bigcup_{j = 1, \ldots, m} V_j$.
For each $j$ we can choose sections
$$
s_{j, 0}, \ldots, s_{j, n_j} \in
\Gamma(f^{-1}(V_j), \mathcal{L}) = f_*\mathcal{L}(V_j) = \mathcal{F}(V_j)
$$
which generate $\mathcal{L}$ over $f^{-1}V_j$ and define an immersion
$$
f^{-1}V_j \longrightarrow \mathbf{P}^{n_j}_{V_j},
$$
see Lemma \ref{lemma-very-ample-finite-type-over-affine}.
Choose $i$ such that there exist sections $e_{j, t} \in \mathcal{E}_i(V_j)$
mapping to $s_{j, t}$ in $\mathcal{F}$ for all $j = 1, \ldots, m$ and
$t = 1, \ldots, n_j$. Then the map $\psi_i$ is surjective
as the sections $f^*e_{j, t}$ have the same image as the sections $s_{j, t}$
which generate $\mathcal{L}|_{f^{-1}V_j}$. Whence we obtain a morphism
$$
r_{\mathcal{L}, \psi_i} : X \longrightarrow \mathbf{P}(\mathcal{E}_i)
$$
over $S$ such that over $V_j$ we have a factorization
$$
f^{-1}V_j \to \mathbf{P}(\mathcal{E}_i)|_{V_j} \to \mathbf{P}^{n_j}_{V_j}
$$
of the immersion given above. It follows that $r_{\mathcal{L}, \psi_i}|_{V_j}$
is an immersion, see Lemma \ref{lemma-immersion-permanence}.
Since $S = \bigcup V_j$ we conclude that $r_{\mathcal{L}, \psi_i}$
is an immersion.
Note that $r_{\mathcal{L}, \psi_i}$ is quasi-compact as
$X \to S$ is quasi-compact and $\mathbf{P}(\mathcal{E}_i) \to S$ is separated
(see Schemes, Lemma \ref{schemes-lemma-quasi-compact-permanence}).
By Lemma \ref{lemma-quasi-compact-immersion} there exists a closed
subscheme $X' \subset \mathbf{P}(\mathcal{E}_i)$ such that $i$ factors
through an open immersion $X \to X'$. Then $X' \to S$ is projective
by definition and we win.
\end{proof}

\end{document}
